\clearpage
\subsection{Original reference: native gas-face ownership and regime}
The working reference is the original three-stage 29/11/4-couple cooler, with $0.40\times0.40\times0.50$ mm P/N legs. At a 300 K reservoir, 100 Pa krypton, 2 mW imposed die heating and 0.595668 W electrical input, the previously solved proxy model gives 219.368939 K at the detector mount. This is a conditional model result, not a measured temperature. Its leg, copper, wire, radiation and gas terms retain the limitations stated above.

The archived nominal connected-gas geometry has 1,840 distinct gas--solid interfaces totaling 903.215 mm$^2$, attached to 429 solid owners. The inventory records each face's location, source lineage and semantic thermal element. Of these faces, 523 belong to patterned conductor nets, 352 to P/N legs, 760 to joints and fillets, and 24 to ceramics. These semantic identifications are not boundary conditions: the local gas-face temperature and exchange law have not yet been mapped to the spatially varying solid models. Eight lead-joint faces retain unresolved thermal destinations. Thus no new gas-field or coupled detector-temperature solution is claimed.

For an engineering rarefaction screen, using an effective krypton collision diameter of 0.36 nm, the hard-sphere mean free path $\lambda=k_{\rm B}T/(\sqrt{2}\pi d^2p)$ is approximately 0.060 mm at 250 K and 100 Pa. It gives $\mathrm{Kn}\approx0.105$ for a 0.570 mm interstage gas-length approximation and $\mathrm{Kn}\approx0.333$ for a 0.180 mm interleg gap. The 0.36 nm kinetic-diameter scale is reported by Kwon and co-workers;\footnote{Y.~H. Kwon et al., \emph{Krypton-Xenon Separation Properties of SAPO-34 Zeolite Materials and Membranes}, https://www.osti.gov/servlets/purl/1356958.} its use as a hard-sphere Kr--Kr diameter is a modeling assumption. These gaps approach or enter the transition regime; ordinary no-slip continuum convection is not established at this pressure. At atmospheric pressure, a separate continuum Rayleigh screen depends strongly on length, orientation and boundary conditions and does not validate the 100 Pa model.
